\begin{abcliste}
\abc Mit $\rho = m/V$ und $n = m/M$ wird aus der idealen Gasgleichung $pV = \frac{m}{M}RT$:
\begin{align}
 \rho &= \frac{pM}{RT}
\end{align}
Sauerstoffgas besteht aus $\mathrm{O_2}$-Molekülen, $M = \MOzF = 2\cdot\MOO = \MOzg$ (FoTaBe S.~228). Mit $T_1 = \TEiF = \TEi$ ist
\begin{align}
 \rho &= \rOF = \frac{\pEi\cdot\MOz}{\ncR\cdot\TEi} = \rO \approx \result{\rOP{0}{2}}
\end{align}

\abc Die Masse ist
\begin{align}
 m &= \mOF = \rO\cdot\VO = \mO \approx \result{\mOP{0}{2}}
\end{align}

\abc Das Volumen bleibt gleich (isochor). Mit $T_2 = \TZwF = \TZw$ ist
\begin{align}
 p_2 &= \pZwF = \pEi\cdot\frac{\TZw}{\TEi} = \pZw \approx \result{\pZwbP{0}{3}}
\end{align}
\end{abcliste}
